\clearpage
\begin{explanationbox}
\textbf{第一眼：定点 $A$ 为什么特别？}

抛物线 $y^2=2px$ 的焦点为 $F(\frac p2,0)$；$A(\frac p2,p)$ 恰在焦点的正上方，同时满足抛物线方程：$p^2=2p\cdot\frac p2$。

从 $A$ 引向抛物线另外两点的两条直线，若斜率分别为 $t$ 与 $-t$，两条直线关于\textbf{经过 $A$ 的水平线}对称。注意：\textbf{对称的是两条直线，不一定是 $M,N$ 两点。}

\medskip
\textbf{第二眼：用一组整点把结论看出来。}
\begin{center}\begin{tikzpicture}[font=\small,>=stealth,line join=round]
\draw[gray!45,rounded corners=3pt] (0,0) rectangle (14.2,9.4);
\node at (7.1,8.8) {两条弦斜率相反，第三条弦的方向固定};
\begin{scope}[shift={(2.4,5.0)},scale=.85]
\draw[->,gray] (-.65,0)--(5.5,0) node[above] {$x$};
\draw[->,gray] (0,-4.8)--(0,3.3) node[left] {$y$};
\draw[teal!65!black,thick,samples=120,domain=-4.5:3.1,variable=\t] plot ({\t*\t/4},\t);
\draw[blue!75!black,very thick] (0,0)--(1,2);
\draw[orange!85!black,very thick] (1,2)--(4,-4);
\draw[green!45!black,very thick] (0,0)--(4,-4);
\fill (1,2) circle(2pt) node[above left=10pt] {$A$};
\fill (0,0) circle(2pt) node[above left=9pt] {$M=O$};
\fill (4,-4) circle(2pt) node[above right=12pt] {$N$};
\fill (1,0) circle(1.8pt) node[below right=10pt] {$F$};
\end{scope}
\node[anchor=west,align=left] at (8.7,5.3)
{$y^2=4x,\quad p=2$\\[12pt]$A(1,2),\quad F(1,0)$\\[12pt]
$M(0,0),\quad N(4,-4)$\\[16pt]
$\textcolor{blue!75!black}{k_{AM}=2},\quad\textcolor{orange!85!black}{k_{AN}=-2}$\\[16pt]
$\textcolor{green!45!black}{k_{MN}=-1}$};
\node[text width=13.2cm,align=center] at (7.1,.65) {坐标轴单位长度相同；斜率相反约束的是两条直线，并非两点坐标。};
\end{tikzpicture}
\end{center}


取 $p=2$，则 $y^2=4x$，$F=(1,0)$，$A=(1,2)$；选取
\[
 M=(0,0),\qquad N=(4,-4).
\]
由图中三组坐标直接回验：$k_{AM}=2$、$k_{AN}=-2$、$k_{MN}=-1$。绿色弦 $MN$ 正好沿着“向右走 $1$、向下走 $1$”的方向。

\medskip
\textbf{第三眼：为什么不用解出两个动点？}

设 $M,N$ 的纵坐标分别为 $u,v$。抛物线上两点的斜率只与\textbf{纵坐标的和}有关：
\[
 \boxed{k_{P(u)P(v)}=\frac{2p}{u+v}}.
\]
题设 $k_{AM}=-k_{AN}$ 恰好使 $u+v=-2p$。因此只求出\textbf{和}，就得到第三条弦的斜率；无须分别求 $M,N$ 的坐标。

\medskip
\textbf{进一步理解：}若把定点换成一般的 $A_0=(\frac{a^2}{2p},a)$（$a\neq0$），同样的斜率对称条件给出 $k_{MN}=-\frac pa$。本题取 $a=p$，所以常数正好化为 $-1$。
\end{explanationbox}
